\section{Results}
\label{sec:results}

\paragraph{Theorem~1 holds to floating-point precision.} In every round of every seed of every NK configuration in the small sweep ($35\!\times\!8 = 280$ measurements), the cosine of $(g_S + g_T)$ against pooled REINFORCE is $1.0000$. Figure~\ref{fig:t1} plots the per-round deviation $1 - \mathrm{cos}$ on a log axis: roughly two thirds of rounds sit at the floating-point floor ($<10^{-15}$) and the rest at $\sim 10^{-7}$, three orders of magnitude tighter than the pre-registered $0.95$ threshold. The identity is sample-wise rather than in expectation, and Corollary~\ref{cor:thm1ar} extends it to autoregressive parameterisations (verified in §\ref{sec:large_scale}). Every claim that follows is grounded in a gradient estimator that is provably and observably equivalent to vanilla REINFORCE up to a partition.

\begin{figure}[t]
\centering
\includegraphics[width=0.6\textwidth]{fig_T1_cosine_histogram.pdf}
\caption{Theorem~\ref{thm:exact} (exactness): per-round deviation $1 - \mathrm{cos}(g_S+g_T,\,\mathrm{pooled\ REINFORCE})$ across all $280$ round-seed pairs in the NK sweep. Maximum deviation is $2.4\!\times\!10^{-7}$.}
\label{fig:t1}
\end{figure}

\paragraph{The Price ratio tracks the analytic optimum.} Across the seven NK configurations, the empirical mean $\rho_t$ correlates with the analytic $\rho^*$ at \textbf{Pearson $r = 0.995$} (Figure~\ref{fig:rho}, right). Both quantities decrease monotonically with the autocorrelation length, in lockstep across nearly two orders of magnitude in $K$: rugged landscapes demand selection-heavy schedules in the low-data regime because each peak found is precious; smooth ones can afford transmission-heavy schedules because exploration is locally informative. The headline best-fitness curve in Figure~\ref{fig:rho} (left) shows that this controller behaviour translates into competitive performance, with PRICE-RL matching AdaLead across the ruggedness range and winning on smooth ($K\!\in\!\{0,1\}$) and very rugged ($K\geq15$) regimes. The $480$-run hyperparameter grid in Appendix~\ref{app:phaseC} confirms that the algorithm is robust to base-LR, inner-step, and support-quantile choice over an order of magnitude on each axis.

\begin{figure}[t]
\centering
\begin{subfigure}[b]{0.49\textwidth}
\includegraphics[width=\textwidth]{fig_E3_best_vs_K.pdf}
\caption{Best fitness reached vs.\ ruggedness $K$.}
\end{subfigure}\hfill
\begin{subfigure}[b]{0.49\textwidth}
\includegraphics[width=\textwidth]{fig_E3_rho_vs_K.pdf}
\caption{Empirical $\rho_t$ vs.\ analytic $\rho^*$ across $K$.}
\end{subfigure}
\caption{NK sweep (E3). H3 strongly supported: $r=0.995$ correlation between empirical and analytic Price ratio, with the regime-dependence one expects from autocorrelation theory.}
\label{fig:rho}
\end{figure}

\paragraph{GB1: tying the active-learning state of the art.} Under the 5-round / 100-query / 5-seed protocol, PRICE-RL achieves top-$1\%$ recall $0.161 \pm 0.047$ vs.\ AdaLead $0.161 \pm 0.050$, PEX $0.031 \pm 0.007$, and random $0.004 \pm 0.001$ (full table in Appendix~\ref{app:phaseB}). The paired-bootstrap test against random gives $p = 0.0000$ and the non-inferiority test against AdaLead gives a gap of $-0.0004$, well within the $1$-std margin of $0.042$. The fixed-mixing ablation collapses to $0.006 \pm 0.003$ --- a $27\times$ regression --- isolating the controller rather than the partition as the active mechanism. The 500-query result is a tie, and we report it as one. In §\ref{sec:large_scale} the lead opens up monotonically with horizon: $0\%$ at $500$ queries, $8\%$ at $4{,}000$, $12.4\%$ at $8{,}000$ with non-overlapping $95\%$ CIs. Figure~\ref{fig:gb1} summarises the small-budget bars; best fitness saturates because the budget exceeds the trivial threshold, so the discriminative metric is top-$1\%$ recall.

\begin{figure}[t]
\centering
\includegraphics[width=0.85\textwidth]{fig_E1_gb1_bars.pdf}
\caption{GB1 (E1). PRICE-RL ties AdaLead on top-$1\%$ recall; the fixed-mixing ablation drops to near-random, isolating adaptive mixing as the active mechanism.}
\label{fig:gb1}
\end{figure}

\paragraph{The Price ratio as an early-warning signal for reward hacking.} Reward hacking is the modal failure of RL-tuned generative biology: the optimisation succeeds against the proxy and silently fails against the true objective. Figure~\ref{fig:e2} plots, for eight independent seeds of the reward-hacking experiment, the trajectory of $\rho_t$ alongside proxy ensemble variance. The Price ratio crosses $0.9$ at round $0$ in every seed --- the policy commits to refining-on-support immediately because the proxy is most informative there --- while proxy variance crosses $95\%$ of its final value only at rounds $9$--$11$, by which point the reward gap has grown past $0.2$. In all $8/8$ seeds the lead of the $\rho_t$ alarm over the proxy-variance alarm is $\geq 9$ rounds. The gradient-level signal beats the behaviour-level signal for a structural reason: $\rho_t$ asks per round whether the policy is consuming new sequence space or refining already-supported variants, and proxy hacking is by construction a refinement-on-artefact event. The §\ref{sec:large_scale} long-horizon study finds a mean lead of $18.6$ rounds across $16$ seeds, and the corruption study shows the diagnostic is robust to surrogate noise of $\sigma_n \in [0, 0.40]$ --- well above the natural noise scale of trained surrogates.

\begin{figure}[t]
\centering
\begin{subfigure}[b]{0.55\textwidth}
\includegraphics[width=\textwidth]{fig_E2_rho_vs_uncertainty.pdf}
\caption{$\rho_t$ vs.\ proxy uncertainty (normalised).}
\end{subfigure}\hfill
\begin{subfigure}[b]{0.42\textwidth}
\includegraphics[width=\textwidth]{fig_E2_reward_gap.pdf}
\caption{Reward gap (proxy $-$ true) over rounds.}
\end{subfigure}
\caption{Reward hacking (E2). The Price ratio fires its $0.9$ alarm $9$--$11$ rounds before proxy variance crosses its $95\%$-of-final threshold on $100\%$ of seeds; the reward gap grows past $0.2$, confirming the policy is exploiting proxy artefacts. A head-to-head against $\delta$-CS (Appendix~\ref{app:phaseB}) shows PRICE-RL detects earlier; $\delta$-CS clips harder.}
\label{fig:e2}
\end{figure}

\paragraph{Cross-domain validation: Trap-5 deceptive landscape.} The Trap-5 landscape \citep{goldberg1989genetic} splits sequences in $\{0,1\}^{30}$ into six 5-bit blocks with per-block fitness $1.0$ at all-ones and a deceptive ramp $f(j) = (4-j)/4$ for $j<5$ ones, so the average gradient points away from the global optimum. Eight rounds, batch $64$, five seeds. Theorem~1 cosine $= 1.0000$ in every round, and PRICE-RL hits the global optimum on $5/5$ seeds; AdaLead lands near $0.95$ (missing the optimum on roughly $1$ in $20$ seeds) and random gets stuck near the deceptive trap. The Price decomposition is the active mechanism: the algorithm escapes deception even though the average gradient is misleading, because the transmission component remains a clean signal of where unexplored mass lies. The §\ref{sec:large_scale} scaling sweep extends this to $\{30,60,120\}\times\{3,5,7,10\}$ Trap-K with PRICE-RL hitting $1.000$ on every seed at every scale, while AdaLead misses up to $20\%$ of the time at $N=120$.

\paragraph{Honest mixed results: cross-protein transfer and long-sequence DMS.} Two studies do \emph{not} support the headline reading. On cross-protein transfer (TEM-1 $\to$ BLAT family), the ranking of cold / full / selection-only transfer modes is target-dependent rather than uniformly favouring the Theorem~3 prescription (Appendix~\ref{app:phaseB}); the clean separation between selection rank-structure and transmission concentration likely requires a non-factorised policy class. On long-sequence DMS (TEM-1 $L=287$, BCL11A enhancer $L=600$), PRICE-RL with uniform-init logits over the full alphabet is not competitive with AdaLead's mutate-from-best locality at small budgets: a fully diffuse factorised policy spreads probability mass over $A^L$ sequences and a few hundred labels are insufficient to concentrate it. The §\ref{sec:large_scale} wildtype-aware initialisation and locality kernel close most of this gap, and the proposal's ESM-2 backbone is the natural plug-in to close the rest.

\paragraph{Aggregate thresholds.} Of the six pre-registered hard thresholds, five pass with large margins: T1 (cosine identity, $\geq 0.95$): $\mathbf{1.0000}$; T3 (Pearson $\geq 0.7$): $r=0.995$; T4 (PRICE vs.\ random, $p<0.10$): $p=0.0000$, gap to AdaLead within $1$ std; T5 ($\rho_t$ leads proxy uncertainty $\geq 2$ rounds, $\geq 80\%$ seeds): $100\%$ of $8$ seeds, $9$--$11$-round lead; T6 (reproducibility from \texttt{scripts/}): all numbers ship from CSV outputs. T2 (PI tight tracking, $|\rho_t-\rho^*|\leq 0.05$ in $100\%$ of seeds): $13.3\%$, below threshold and structurally so (§\ref{sec:discussion}). Full per-experiment tables and the threshold audit are in Appendix~\ref{app:phaseB}--\ref{app:audit}.
